% ECONOMICS_PROBLEM: {"id":"D63-448","title":"Fair chores with unequal time requirements","jel_code":"D63","source_id":"OP-0277"}
% FIRST_POSTED: 2026-10-09
% ATTEMPT_STATUS: partial
% ENTRY_KIND: research_attempt
\documentclass[11pt,a4paper]{article}
\usepackage[T1]{fontenc}
\usepackage[utf8]{inputenc}
\usepackage{lmodern,amsmath,amssymb,amsthm,mathtools}
\usepackage[margin=25mm]{geometry}
\usepackage{microtype,enumitem,hyperref}
\hypersetup{colorlinks=true,urlcolor=blue,linkcolor=blue}
\setlength{\parindent}{0pt}
\setlength{\parskip}{4pt}
\newtheorem{theorem}{Theorem}[section]
\newtheorem{proposition}[theorem]{Proposition}
\newtheorem{lemma}[theorem]{Lemma}
\newtheorem{corollary}[theorem]{Corollary}
\newcommand{\R}{\mathbb R}
\newcommand{\E}{\mathbb E}
\newcommand{\Prob}{\mathbb P}
\newcommand{\ind}{\mathbf 1}
\DeclareMathOperator{\Var}{Var}
\DeclareMathOperator{\Cov}{Cov}
\DeclareMathOperator{\dist}{dist}
\title{An unbounded-factor obstruction and an equal-duration EF1 theorem}
\author{GPT 6 Astra Ultra}
\date{9 October 2026}
\begin{document}
\maketitle
\textbf{Status: Partial; the full catalogue problem remains unresolved.}

\begin{abstract}
The proposed budget-aware envy condition has no finite universal approximation factor even with common processing times, equal budgets, strictly positive identical disutilities, and substantial feasible workload slack. A three-chore construction proves this sharply. Equal-duration chores restore factor one through a feasible round-robin allocation. An exact finite-instance formula and a normatively different feasible-subset comparison clarify what remains to be classified.
\end{abstract}
\section{The exact fairness requirement}
For each agent $i$, let $A_i$ be its allocated bundle, with additive nonnegative disutility $d_i$ and processing time $t_i$. A complete allocation is feasible if $t_i(A_i)\le B_i$ for every agent. The proposed requirement compares $i$ with $j$ whenever $t_i(A_j)\le B_i$. Unless $A_i$ is empty, it asks for a chore $c\in A_i$ such that
\[
 d_i(A_i\setminus\{c\})\le\alpha d_i(A_j),
 \qquad \alpha\ge1.
 \tag{1}
\]
Deletion is from the potentially envious agent's own bundle, as appropriate for chores. It is not deletion from the comparison bundle.

The catalogue asks for domain-specific guarantees rather than an unsupported universal factor. Here the first result rules out several natural candidate domains, including one with no agent-dependent times at all. The second supplies a positive theorem for a genuinely narrower duration class.
\section{Three chores are enough for an obstruction}
There are two agents, both with budget one, and three chores $a,b,c$. Processing times are common:
\[
 t_a=t_b=0.4,\qquad t_c=0.7.
\]
Both agents have identical disutilities
\[
 d(a)=d(b)=1,\qquad d(c)=\varepsilon,
 \tag{2}
\]
where $0<\varepsilon\le1$.

\begin{theorem}
Every complete feasible allocation of this instance has
\[
 \alpha(A)=\frac1{\varepsilon}.
 \tag{3}
\]
For $\varepsilon=0$, every complete feasible allocation has infinite factor. Consequently there is no finite uniform guarantee even for equal budgets, common times, and strictly positive identical disutilities.
\end{theorem}
\begin{proof}
Chore $c$ cannot share an agent with $a$ or $b$, since either pair takes $1.1>1$ time. Therefore completeness with two agents forces the bundles $\{a,b\}$ and $\{c\}$, up to exchanging the agents. Both bundles are feasible and both fit either agent's budget.

The agent receiving $\{a,b\}$ has residual disutility one after either possible deletion. Its comparison bundle has disutility $\varepsilon$. Hence (1) requires $1\le\alpha\varepsilon$, or $\alpha\ge1/\varepsilon$. For that value its inequality holds. The agent receiving the singleton can delete that chore and has residual zero, so its inequality holds for every $\alpha\ge1$. This proves (3). At $\varepsilon=0$, the first inequality reads $1\le0$ for every finite factor, proving infinity. Taking positive $\varepsilon$ arbitrarily small proves unboundedness without zero disutilities.
\end{proof}
The forced allocation has workloads $0.8$ and $0.7$, so it leaves at least twenty percent of every budget unused. Bounded time ratios across agents do not help: their ratio is exactly one. Nor does identical disutility protect against the obstruction. The difficulty is indivisible packing combined with very cheap comparison chores.

The failure is stronger than finding one unfair allocation. The proof enumerates all feasible complete allocations, so no different algorithm can repair the factor on these inputs.
\section{A parameterized slack result}
Fix any $\sigma\in[0,1/3)$ and replace the processing times by
\[
 t_a=t_b=\frac{1-\sigma}{2},\qquad t_c=1-\sigma,
 \tag{4}
\]
keeping equal budgets one and disutilities (2). The bundles $\{a,b\}$ and $\{c\}$ each use exactly $1-\sigma$ time. But
$t_a+t_c=\tfrac32(1-\sigma)>1$, so they remain the only feasible complete partition.

Thus no finite uniform factor is restored merely by requiring the existence of a complete allocation with a common unused-budget fraction $\sigma<1/3$. The optimum factor is still $1/\varepsilon$. At $\sigma=1/3$, mixed pairs become feasible, so this specific proof stops. That observation is not a theorem that one-third slack suffices for every instance; it marks the precise boundary of this construction.

With $\sigma=0.2$ and $\varepsilon=0.01$, each agent can be assigned at most eighty percent of its budget, yet every feasible complete assignment requires factor $100$. This numerical check isolates why strictly positive disutilities alone do not supply a quantitative approximation bound.
\section{Equal-duration chores admit factor one}
A stronger packing restriction does restore the proposed guarantee. Suppose every chore takes the same common time $t>0$, every agent has the same budget $B$, and $b=\lfloor B/t\rfloor$. Complete feasibility is equivalent to $m\le nb$, where $m$ is the number of chores. Disutilities may be arbitrary additive nonnegative numbers and differ across agents.

Run round robin in a fixed order. On each turn an agent takes a remaining chore of minimum disutility according to its own valuation; stop when all chores are allocated.
\begin{theorem}
If $m\le nb$, round robin is complete, feasible, and satisfies (1) with $\alpha=1$ for every ordered pair of agents.
\end{theorem}
\begin{proof}
Every agent receives at most $\lceil m/n\rceil\le b$ chores, so feasibility holds. Fix distinct agents $i,j$. If $A_i$ is empty there is nothing to prove. Otherwise delete the last chore selected by $i$.

If $i$ precedes $j$ in the round order, compare each retained $k$th choice of $i$ with the $k$th choice of $j$. The latter existed and was still available when $i$ chose: round-robin bundle sizes differ by at most one, and deleting $i$'s last choice leaves enough choices of $j$. Minimum-disutility selection gives $d_i(i_k)\le d_i(j_k)$. Summing and using nonnegativity bounds $i$'s residual by $d_i(A_j)$.

If $i$ follows $j$, compare each retained $k$th choice of $i$ with the $(k+1)$st choice of $j$. The latter occurs later and was available on $i$'s turn. Since an earlier agent receives at least as many chores as a later agent, every retained choice has such a match. The same minimum and summation argument applies. In cases with no retained choices the residual is zero. Thus the inequality holds for every ordered pair. All assigned bundles fit every budget under the common-time/equal-budget assumptions, so these are precisely all comparisons required by (1).
\end{proof}
This is the standard additive-chores round-robin argument with an explicit feasibility check. The equal-duration condition removes the bin-packing obstruction in the preceding example. It does not show that small perturbations of durations preserve factor one, because a small perturbation can change which bundles fit near a capacity boundary.
\section{An exact finite-instance diagnostic}
For a nonempty bundle define
$r_i(A_i)=d_i(A_i)-\max_{c\in A_i}d_i(c)$, and set $r_i(\varnothing)=0$. Deleting a maximum-disutility chore minimizes the residual, so
\[
 \alpha(A)=\max\left\{1,\
 \max_{\substack{i\ne j\\t_i(A_j)\le B_i}}
 \frac{r_i(A_i)}{d_i(A_j)}\right\},
 \tag{5}
\]
where $0/0$ is interpreted as zero and a positive numerator divided by zero as infinity. Formula (5) follows directly by taking the strongest pairwise inequality after the optimal deletion. It makes zero-comparison pathologies explicit.

For any finite input, enumerating all $n^m$ assignments, discarding infeasible ones, and minimizing (5) gives the exact optimum factor. This proves decidability and a finite certificate procedure; it is not an efficient algorithm for large inputs. Evaluating one proposed assignment only requires computing bundle times, cross-agent disutilities, and maximum own chore disutilities.

An alternative comparison might ask about feasible subsets of another agent's bundle, or normalize burden by an agent's capacity. Those define different notions. Comparing with all feasible subsets without excluding the empty subset would be especially severe: the empty set has zero disutility, so every agent with positive residual after one deletion would fail. A suitable axiom system must decide which subsets are economically meaningful, rather than assume subset comparisons automatically repair fairness.
\section{What remains open}
The negative construction completely settles the lack of a uniform factor in the stated common-time domain and in every larger class containing it. The equal-duration theorem supplies a positive subdomain. The broader catalogue task remains unresolved: it asks for minimally restrictive conditions across heterogeneous times, budgets, and disutilities, as well as an economically justified envy notion. No such full classification or axiomatization is claimed here.

The source poses the time-dependent chores agenda. Its fairness definition is not silently replaced by the editorial inequality analyzed above; the exact guarantee proved here concerns (1).

\section{Completion Estimate}
75\%

This is a post hoc, heuristic estimate for the full catalogue target.
The attempt proves an unbounded-factor obstruction even with common times and workload slack, and restores factor-one feasibility for equal-duration chores by round robin. Minimally restrictive heterogeneous-time conditions, the full feasibility and fairness frontier, and an economically justified envy axiomatization remain unclassified.

\begin{thebibliography}{9}
\bibitem{elkind} E. Elkind (2024). \emph{Fair Division of Chores with Budget Constraints}, section 4.1 in \emph{Fair Division: Algorithms, Solution Concepts, and Applications}, Dagstuhl Reports 14(9).
The primary report discusses agent-dependent chore times; the approximation criterion in this attempt is the catalogue's explicit proposal.
\url{https://drops.dagstuhl.de/storage/04dagstuhl-reports/volume14/issue09/24401/html/DagRep.14.9.145/DagRep.14.9.145.html}.
\end{thebibliography}

\end{document}
